% GENERATED FILE. Edit canonical lessons, then run npm run generate:books.
% canonical-source-hash: fnv1a64:9d0ea95d61225c10
% canonical-lessons: ES-C761-serie, ES-C761-videojuego, ES-C761-cojin, ES-C761-bostezar, ES-C761-grabar

\chapter{Tardes en el sofá}
\label{ch:es-tardes-en-el-sofa}
\hlchaptermodality{761}

\begin{chapteropening}
\textbf{By the end of this chapter:} \emph{I can describe a lazy afternoon at home: a series, video games, cushions and yawns.}

Complete the last lesson of chapter 761: I can describe a lazy afternoon at home: a series, video games, cushions and yawns.
\end{chapteropening}

\section[la serie]{la serie — the series (TV)}

\label{lesson:ES-C761-serie}



\begin{quote}
Before the new one: say the Spanish for the premiere, then the Spanish for to play.
\end{quote}

\subsection*{You'll want to know: la serie}
\textbf{la serie} — \textquotedblleft{}the series (TV)\textquotedblright{}.

\begin{quote}
Estoy viendo una serie muy buena.
\end{quote}

\begin{quote}
\emph{I am watching a really good series.}
\end{quote}

Same form in the plural: \emph{las series}.

\begin{cousinweb}
\emph{serie} comes from Latin \emph{seriēs} \textquotedblleft{}a row, a chain\textquotedblright{} (EtymDB), from \emph{serere} \textquotedblleft{}to join together\textquotedblright{}.
\end{cousinweb}

\begin{grammarlens}[title={What you've built}]
One more word for telling, explaining and describing at B1.
\end{grammarlens}

\subsection*{Your turn}
\begin{itemize}
\raggedright
  \item \emph{Say it:} \emph{la serie}
  \item \emph{Say it:} \emph{la serie}, once more
  \item \emph{Say it:} a sentence of your own with \emph{la serie}
  \item \emph{Recall:} say the Spanish for the premiere, then the Spanish for to play, then say \emph{la serie} again
\end{itemize}

\begin{tcolorbox}[breakable,colback=teal!4,colframe=teal!35!black,title={Before you move on}]
What does la serie mean? (\textquotedblleft{}The series (TV)\textquotedblright{}.) Say it once more.
\end{tcolorbox}
\section[el videojuego]{el videojuego — the video game}

\label{lesson:ES-C761-videojuego}



\begin{quote}
Before the new one: say the Spanish for to play, then the Spanish for the series.
\end{quote}

\subsection*{You'll want to know: el videojuego}
\textbf{el videojuego} — \textquotedblleft{}the video game\textquotedblright{}.

\begin{quote}
Mi hijo juega a los videojuegos.
\end{quote}

\begin{quote}
\emph{My son plays video games.}
\end{quote}

\emph{Jugar a los videojuegos}.

\begin{cousinweb}
\emph{videojuego} is built on \emph{vídeo}, from Latin \emph{videō} \textquotedblleft{}I see\textquotedblright{}, and \emph{juego} \textquotedblleft{}game\textquotedblright{}, from Latin \emph{iocus} \textquotedblleft{}jest\textquotedblright{} (EtymDB); it follows English \emph{video game}.
\end{cousinweb}

\begin{grammarlens}[title={What you've built}]
One more word for telling, explaining and describing at B1.
\end{grammarlens}

\subsection*{Your turn}
\begin{itemize}
\raggedright
  \item \emph{Say it:} \emph{el videojuego}
  \item \emph{Say it:} \emph{el videojuego}, once more
  \item \emph{Say it:} a sentence of your own with \emph{el videojuego}
  \item \emph{Recall:} say the Spanish for to play, then the Spanish for the series, then say \emph{el videojuego} again
\end{itemize}

\begin{tcolorbox}[breakable,colback=teal!4,colframe=teal!35!black,title={Before you move on}]
What does el videojuego mean? (\textquotedblleft{}The video game\textquotedblright{}.) Say it once more.
\end{tcolorbox}
\section[el cojín]{el cojín — the cushion}

\label{lesson:ES-C761-cojin}



\begin{quote}
Before the new one: say the Spanish for the series, then the Spanish for the video game.
\end{quote}

\subsection*{You'll want to know: el cojín}
\textbf{el cojín} — \textquotedblleft{}the cushion\textquotedblright{}.

\begin{quote}
El sofá tiene muchos cojines.
\end{quote}

\begin{quote}
\emph{The sofa has lots of cushions.}
\end{quote}

Plural \emph{cojines}.

\begin{cousinweb}
\emph{cojín} comes from Old Spanish \emph{coxín}, from a spoken Latin \emph{coxīnum} (EtymDB), built on \emph{coxa} \textquotedblleft{}hip\textquotedblright{}: something to rest the hip on. Many dictionaries say it came in through French \emph{coussin}.
\end{cousinweb}

\begin{grammarlens}[title={What you've built}]
One more word for telling, explaining and describing at B1.
\end{grammarlens}

\subsection*{Your turn}
\begin{itemize}
\raggedright
  \item \emph{Say it:} \emph{el cojín}
  \item \emph{Say it:} \emph{el cojín}, once more
  \item \emph{Say it:} a sentence of your own with \emph{el cojín}
  \item \emph{Recall:} say the Spanish for the series, then the Spanish for the video game, then say \emph{el cojín} again
\end{itemize}

\begin{tcolorbox}[breakable,colback=teal!4,colframe=teal!35!black,title={Before you move on}]
What does el cojín mean? (\textquotedblleft{}The cushion\textquotedblright{}.) Say it once more.
\end{tcolorbox}
\section[bostezar]{bostezar — to yawn}

\label{lesson:ES-C761-bostezar}



\begin{quote}
Before the new one: say the Spanish for the video game, then the Spanish for the cushion.
\end{quote}

\subsection*{You'll want to know: bostezar}
\textbf{bostezar} — \textquotedblleft{}to yawn\textquotedblright{}.

\begin{quote}
Bostezo mucho cuando tengo sueño.
\end{quote}

\begin{quote}
\emph{I yawn a lot when I am sleepy.}
\end{quote}

\emph{El bostezo} — the yawn. Spelling: \emph{bostecé}.

\begin{cousinweb}
\emph{bostezar} is usually traced to Latin \emph{ōscitāre} \textquotedblleft{}to gape, to yawn\textquotedblright{} (EtymDB); how it gained its \emph{b-} is not clear.
\end{cousinweb}

\begin{grammarlens}[title={What you've built}]
One more word for telling, explaining and describing at B1.
\end{grammarlens}

\subsection*{Your turn}
\begin{itemize}
\raggedright
  \item \emph{Say it:} \emph{bostezar}
  \item \emph{Say it:} \emph{bostezar}, once more
  \item \emph{Say it:} a sentence of your own with \emph{bostezar}
  \item \emph{Recall:} say the Spanish for the video game, then the Spanish for the cushion, then say \emph{bostezar} again
\end{itemize}

\begin{tcolorbox}[breakable,colback=teal!4,colframe=teal!35!black,title={Before you move on}]
What does bostezar mean? (\textquotedblleft{}To yawn\textquotedblright{}.) Say it once more.
\end{tcolorbox}
\section[grabar]{grabar — to record; to engrave}

\label{lesson:ES-C761-grabar}



\begin{quote}
Before the new one: say the Spanish for the cushion, then the Spanish for to yawn.
\end{quote}

\subsection*{You'll want to know: grabar}
\textbf{grabar} — \textquotedblleft{}to record; to engrave\textquotedblright{}.

\begin{quote}
Grabamos el concierto con el móvil.
\end{quote}

\begin{quote}
\emph{We recorded the concert on our phones.}
\end{quote}

For sound and video; also to engrave.

\begin{cousinweb}
\emph{grabar} was borrowed from French \emph{graver} \textquotedblleft{}to engrave\textquotedblright{} (EtymDB), which comes from Frankish \emph{graban} \textquotedblleft{}to dig\textquotedblright{}; a recording is sound cut into a surface.
\end{cousinweb}

\begin{grammarlens}[title={What you've built}]
One more word for telling, explaining and describing at B1.
\end{grammarlens}

\subsection*{Your turn}
\begin{itemize}
\raggedright
  \item \emph{Say it:} \emph{grabar}
  \item \emph{Say it:} \emph{grabar}, once more
  \item \emph{Say it:} a sentence of your own with \emph{grabar}
  \item \emph{Recall:} say the Spanish for the cushion, then the Spanish for to yawn, then say \emph{grabar} again
\end{itemize}

\begin{tcolorbox}[breakable,colback=teal!4,colframe=teal!35!black,title={Before you move on}]
What does grabar mean? (\textquotedblleft{}To record; to engrave\textquotedblright{}.) Say it once more.
\end{tcolorbox}
